\subsection{Quelques analyses de sensibilité}
\label{sec:sensibilite_sans_retrait}

\subsubsection{Transmission en ligne directe avec patrimoine important}
\label{sec:sensibilite_directe}

\begin{figure}[h!]
    \centering
    \includegraphics[width=\textwidth]{cas_4.png}
    \caption{Différence relative d'héritage net pour un investissement de 1 000 000 euros avant 70 ans, avec un patrimoine hors enveloppes de 1 000 000 euros, en cas de transmission directe.}
    \label{fig:sim4}
\end{figure}

Dans ce scénario d'investissement d'un million d'euros en assurance-vie ou en compte-titres ordinaire, adossé à un patrimoine hors enveloppes de même ampleur, l'avantage relatif du CTO apparaît marqué et les conclusions sont similaires à celles faites précédemment pour un investissement plus faible. Le panneau de gauche de la figure \ref{fig:sim4} montre que l'effet cumulatif des frais de gestion, même modérés, érode rapidement la performance successorale de l'assurance-vie à mesure que la durée de détention s'allonge. Le panneau de droite illustre que, pour des frais fixes de 0,5\%, l'assurance-vie ne conserve un léger avantage que dans des situations défavorables à l'épargnant (rendements faibles ou décès précoce). Dès que le rendement croît et que l'horizon s'étend, le CTO reprend l'avantage grâce à la neutralisation des plus-values latentes au décès. L'abattement spécifique de l'assurance-vie devient marginal à ces niveaux de patrimoine, et ne suffit pas à compenser le poids des frais et des prélèvements sociaux. En ligne directe, et pour des masses patrimoniales élevées, le CTO se révèle ainsi plus efficient que l'assurance-vie dans la majorité des configurations simulées. Enfin, le passage dans les tranches supérieures du barème général n'altère pas fondamentalement ces résultats : l'écart relatif entre les deux enveloppes reste dominé par l'effet des frais et des prélèvements sociaux propres à l'assurance-vie.

\subsubsection{Transmission en ligne indirecte}
\label{sec:sensibilite_indirecte}

Cette section analyse deux configurations de transmission en ligne indirecte (neveu ou nièce), construites à partir des figures \ref{fig:sim5} et \ref{fig:sim6}. Dans les deux cas, l'investissement étudié s'élève à 1 000 000 euros et le patrimoine hors enveloppes à 1 000 000 euros. Les cartes présentent la différence relative d'héritage net
\begin{equation}
    \Delta H_{rel} = \frac{H_{AV} - H_{CTO}}{CTO + \text{autres biens}},
\end{equation}
en fonction de la durée de détention et, selon le panneau, des frais de gestion de l'assurance-vie ou du rendement annuel. Les valeurs positives indiquent un avantage de l'assurance-vie, les valeurs négatives un avantage du compte-titres ordinaire (CTO).

\begin{figure}[h!]
    \centering
    \includegraphics[width=\textwidth]{cas_5.png}
    \caption{Ligne indirecte (neveu ou nièce), patrimoine hors enveloppes 1 000 000 euros, investissement 1 000 000 euros. Panneau gauche : impact des frais de gestion de l'assurance-vie pour un rendement fixé à 4,5\% l'an. Panneau droit : impact du rendement pour des frais fixes de 0,5\% l'an.}
    \label{fig:sim5}
\end{figure}

\begin{figure}[h!]
    \centering
    \includegraphics[width=\textwidth]{cas_6.png}
    \caption{Ligne indirecte (neveu ou nièce), même configuration mais avec des frais de gestion de 2,0\% l'an pour l'assurance-vie dans le panneau de droite et un rendement fixé à 4,5\% l'an dans le panneau de gauche.}
    \label{fig:sim6}
\end{figure}

Avec des frais de gestion faibles, l'assurance-vie conserve une zone étendue de supériorité. Cet avantage provient de la conjonction de l'abattement spécifique par bénéficiaire et des taux applicables à l'assurance-vie, plus modérés que le taux forfaitaire de 55\% en ligne indirecte. L'augmentation du rendement réduit cet écart lorsque la durée de détention s'allonge, toutes choses égales par ailleurs puisque dans ce cas, la purge des plus values lors de la transmission du CTO compense en partie le taux d'imposition plus important.

Lorsque les frais atteignent des niveaux plus représentatifs des contrats français, l'avantage relatif de l'assurance-vie se réduit sensiblement. À horizons longs, la combinaison de la charge de frais cumulée et des prélèvements sociaux sur les gains peut annuler l'avantage théorique conféré par l'abattement et le barème spécifiques, voire aboutir à une supériorité du CTO. Le rôle de la neutralisation des plus-values latentes au décès dans le CTO devient alors déterminant.
